% \subsection{SynchronousCommandExecutorImplementation} 
% \begin{algorithm}[H]
% \caption{SynchronousCommandExecutorImplementation (\emph{Class}) One of Multiple Selectable Executors (Sync, Async etc) via Configuration Properties}
% \begin{algorithmic}[1]
% \Procedure{Execute}{command : Command\textlangle REQ\textrangle}
%     \State \textbf{Input:} command of type \textlangle REQ\textrangle
%     \State \textbf{Output:} Supplier of type \textlangle RES\textrangle
%     \State \textbf{Description:} Executes the command synchronously by invoking all middlewares and routing to the appropriate handler.
%     \State
%     \State \textit{Note:} This implementation is conditionally activated based on application configuration
%     (e.g., \texttt{command.executor = sync, async, disruptor etc}), 
%     allowing multiple executor implementations to coexist and be selected at runtime.
%     \State
%     \ForAll{middleware \textbf{in} middlewares}
%         \State middleware.Invoke(command)
%     \EndFor
%     \State handler $\gets$ router.Route(command)
%     \State \textbf{Return:} a Supplier that, when invoked, executes handler.Handle(command) and returns the result
% \EndProcedure
% \end{algorithmic}
% \end{algorithm}

\subsection{SynchronousCommandExecutorImplementation}

The \texttt{SynchronousCommandExecutorImplementation} represents one of multiple interchangeable executor variants 
(e.g., \texttt{Sync}, \texttt{Async}, \texttt{Disruptor-based}) selectable via configuration.  
This executor performs \emph{synchronous, blocking} command execution by sequentially invoking middleware and dispatching the command to its handler.

At runtime, the executor is conditionally activated based on configuration properties such as:

\[
\texttt{command.executor} \in \{ \texttt{sync}, \texttt{async}, \texttt{disruptor}, \ldots \}
\]

Formally, for a command $c_{\tau_P}$, the executor defines:
\[
E_{\text{sync}}(c_{\tau_P}) = \lambda() \rightarrow h_{\tau_H}(M(c_{\tau_P}))
\]
where $M$ represents the ordered composition of middlewares and $h_{\tau_H}$ is the routed handler.

\begin{algorithm}[H]
\caption{SynchronousCommandExecutorImplementation (\emph{Class}) \\ 
One of Multiple Selectable Executors (Sync, Async, etc.) via Configuration Properties}
\begin{algorithmic}[1]

\Procedure{Execute}{command : Command\textlangle REQ\textrangle}
    \State \textbf{Input:} command of type \textlangle REQ\textrangle
    \State \textbf{Output:} Supplier of type \textlangle RES\textrangle
    \State \textbf{Description:} Executes the command synchronously by invoking all middlewares and routing to the appropriate handler.
    \State
    \State \textit{Note:} This implementation is conditionally activated based on application configuration
    (e.g., \texttt{command.executor = sync, async, disruptor, etc.}),
    allowing multiple executor implementations to coexist and be selected at runtime.
    \State
    \ForAll{middleware \textbf{in} middlewares}
        \State middleware.Invoke(command)
        \Comment{Perform validation, logging, or transformation}
    \EndFor
    \State handler $\gets$ router.Route(command)
        \Comment{Determine the appropriate handler for this command type}
    \State \textbf{Return:} a Supplier that, when invoked, executes \texttt{handler.Handle(command)} and returns the result
\EndProcedure

\end{algorithmic}
\end{algorithm}

This synchronous executor enforces immediate evaluation — the handler logic is executed in the same thread of invocation.
It is best suited for low-latency, deterministic workflows where blocking behavior is acceptable.  
Alternative implementations (e.g., \texttt{AsyncCommandExecutor}, \texttt{DisruptorCommandExecutor}) can extend this pattern to achieve non-blocking or event-driven semantics.